\documentclass[11pt]{article}
\usepackage[margin=1in]{geometry}
\usepackage[T1]{fontenc}
\usepackage{lmodern,amsmath,amssymb,amsthm,mathtools,hyperref}
\setlength{\emergencystretch}{3em}
\newtheorem{theorem}{Theorem}[section]
\newtheorem{lemma}[theorem]{Lemma}
\newtheorem{corollary}[theorem]{Corollary}
\newtheorem{proposition}[theorem]{Proposition}
\newcommand{\E}{\mathbb E}
\newcommand{\tr}{\operatorname{tr}}
\title{An exact R\'enyi-order transition\\for four pairwise-independent bits}
\author{Joshua Mark\thanks{ORCID: \href{https://orcid.org/0009-0003-5905-2382}{0009-0003-5905-2382}. AI tools assisted with derivation, independent review, computation and exposition.}}
\date{20 November 2026 --- research draft; priority not established}
\begin{document}\maketitle
\section{An exact entropy-order transition on four bits}
\label{sec:renyi-phase}
An entropy-order continuation method would vary the objective in the
hope that the desired Hadamard distribution becomes easier to reach.
The smallest nontrivial moment polytope already has a complete,
explicit transition with persistent competing local optima. The
following analysis is self-contained; it does not establish novelty
or an all-order existence theorem.

Let $\mathcal P_4$ be the laws on four sign bits with zero means and
zero distinct pair correlations. For $0<\alpha<\infty$, $\alpha\ne1$,
define
\[
 H_\alpha(p)=\frac{\log\sum_xp_x^\alpha}{1-\alpha},
\]
with $H_1$ the Shannon entropy and all logarithms natural.
Define two vertex families: $\mathcal U$ consists of the ten uniform
laws obtained by fixing one triple parity or the quadruple parity;
$\mathcal W$ consists of the sixteen sign translations of
\[
 w_x=\begin{cases}
 1/6,&k(x)=0,\\
 1/12,&k(x)\in\{2,3\},\\
 0,&k(x)\in\{1,4\},
 \end{cases}
\]
where $k$ counts negative coordinates.

\begin{theorem}[Complete four-bit entropy diagram]\label{thm:renyi-phase}
There is a unique $\alpha_*\in(6,7)$ satisfying
\[
 (4/3)^{\alpha_*}+10(2/3)^{\alpha_*}=8.
\]
For every finite $\alpha>0$, the global minimizers of $H_\alpha$ on
$\mathcal P_4$ are exactly $\mathcal U$ when $\alpha<\alpha_*$,
exactly $\mathcal W$ when $\alpha>\alpha_*$, and their union at
$\alpha=\alpha_*$. Explicitly, the minimum is the smaller of
\[
 \log8,\qquad
 \frac{\log(2^\alpha+10)-\alpha\log12}{1-\alpha},
\]
with the second expression interpreted by continuity at one.
Every law in both families remains a strict local minimum for every
finite $\alpha>1$. Thus varying the entropy order does not force
the competing local minima to disappear at the global transition.
\end{theorem}

\begin{proof}
Fourier expansion writes every feasible law uniquely as
\[
 16p_x=1+\sum_{i=1}^4a_i\chi_{[4]\setminus\{i\}}(x)
                   +a_5\chi_{[4]}(x).
\]
The five displayed characters take exactly the sixteen sign patterns
$v\in\{\pm1\}^5$ with $\prod_i v_i=1$. Hence the parameter polytope is
\[
 P=\{a\in\mathbb R^5:1+v\cdot a\ge0\ \text{for all such }v\}.
\]
We first prove the exact vertex description
\[
 P=\operatorname{conv}\left(
 \{\pm e_i:1\le i\le5\}
 \cup\{s/3:s\in\{\pm1\}^5,\ \prod_i s_i=1\}\right).
\]
All listed points satisfy the inequalities. Conversely put
$b_i=|a_i|$ and $m=\min_i b_i$. If some coordinate is zero or the
product of coordinate signs is negative, an admissible $v$ gives
$v\cdot a=-\sum_i b_i$, so $\sum_i b_i\le1$ and $a$ lies in the
cross-polytope generated by the signed coordinate vectors.
Otherwise let $s_i=\operatorname{sign}(a_i)$, with $\prod s_i=1$.
The least admissible dot product is $-\sum_i b_i+2m$, so
$\sum_i b_i-2m\le1$. The representation
\[
 a=(3m)(s/3)+\sum_i(b_i-m)s_i e_i
\]
has nonnegative coefficients of total $\sum_i b_i-2m\le1$;
assign any remaining mass to zero, which is in the cross-polytope.
This proves the hull identity. The signed coordinate vectors are
unique extrema of coordinate functionals. Each $s/3$ uniquely maximizes
$s\cdot a$ among the listed points, attaining $5/3$ rather than at
most one elsewhere. Thus all twenty-six points are vertices.
They yield precisely $\mathcal U$ and $\mathcal W$ above.

For $\alpha>1$, minimizing entropy means maximizing the strictly
convex function $\sum p^\alpha$; for $0<\alpha<1$ it means minimizing
that strictly concave function. At one Shannon entropy is strictly
concave. In every case a nonvertex cannot minimize, so it suffices
to compare the two weight multisets. For $\alpha\ne1$ their power sums
differ in sign according to
\[
 \phi(\alpha)=(4/3)^\alpha+10(2/3)^\alpha-8.
\]
This function is strictly convex. It has a root at one, and exact
rational evaluation gives $\phi(6)<0<\phi(7)$. There is exactly
one other root, $\alpha_*\in(6,7)$; the signs are positive before
one, negative between one and $\alpha_*$, and positive afterwards.
Accounting for the sign of $1-\alpha$ proves the stated comparison.
At one the weighted entropy is
$\log12-(\log2)/6>\log8$. Strict convexity or concavity also excludes
nonvertex mixtures at the crossing.

It remains to verify local persistence, not merely stationarity.
Fix finite $\alpha>1$. For a weighted vertex put
$u=(1/6)^{\alpha-1}$ and $b=(1/12)^{\alpha-1}$. The degree-two function
\[
 q(k)=b+\frac{u-b}{6}(k-2)(k-3)
\]
equals $w_x^{\alpha-1}$ on its positive support and is strictly
positive elsewhere. Its expectation is fixed by the moment constraints.
It therefore strictly exposes the gradient at $w$, unless a feasible
law has the same support. There is no other feasible law on that
support: a second one would allow a small continuation from $w$ in
the opposite direction and contradict that $w$ is a vertex.
For a uniform vertex the same argument uses the constant polynomial
equal to its nonzero gradient value.

For clarity, strict linear exposure suffices here because the feasible
set is a finite polytope. The negative gradient slope divided by the
distance to each other vertex has a positive minimum. Convex
decomposition and the triangle inequality give a uniform negative
linear bound in every feasible direction. Differentiability of
$\sum p^\alpha$ for $\alpha>1$ makes its remainder $o(\|p-w\|)$,
proving strict local maximality of the power sum and hence strict
local minimality of entropy. This argument includes $1<\alpha<2$,
without requiring a bounded Hessian at zero coordinates.
\end{proof}

The order-one root of $\phi$ reflects normalization and is not an
entropy transition. The genuine transition lies between six and seven.
As $\alpha$ tends to infinity the weighted optimum approaches
$\log6$, whereas the collision and Shannon optima are $\log8$.
The local-persistence statement explicitly excludes the endpoint
$\alpha=\infty$.

This example is a calibrated obstruction to using entropy-order
continuation alone as a global attainment argument. It does not
exclude nonlocal moves or an independently justified global method.
It also illustrates a useful geometric reduction: the four-bit
moment polytope is the polar of a five-dimensional parity half-cube,
and its entire optimization diagram follows from its explicit hull.

\end{document}
